% Shared v5 scheduler chapter, included by both v4-layout user guides.
\section{The three \code{parfeval} schedulers}\label{sec:schedulers}

\subsection{Execution framework versus scheduler}
The \emph{framework} determines how curve traces execute: serial, synchronous
\code{parfor}, or asynchronous \code{parfeval}. The \emph{scheduler} determines
which eligible trace the \code{parfeval} queue launches next. Therefore
\code{HEBCPOLICY} affects only \code{runVBook_hybrid_parfeval.m}; it does not
alter serial or \code{parfor} mathematics. All policies share the branch tracer,
corrector, Pad\'e calculation, deduplication, residual tests, and completion rule.

A \emph{trace} follows one continuation curve from one known solution for one
equation. It can return zero, one, or several previously unknown solutions.
Thus trace count measures completed work, not solution count. Trace-to-90\% is
the completion index at which accumulated discoveries first reach 90\% of that
run's final solution total.

\subsection{Shared queue state and completeness}
The queue operates on pairs $(s,e)$ of known solution and equation.
\code{VBook(s,e)==0} marks an uncovered pair. Only a non-busy equation can be
dispatched, so two workers never race to stamp the same equation column.
Each equation owns a monotonic solution-index cursor that advances past covered
rows and later reaches newly appended solutions. This replaces v4's duplicated
pending lists: auxiliary scheduler state is $O(n_{\rm eq})$, while the necessary
coverage matrix remains $O(n_s n_{\rm eq})$.

\subsection{\code{scan}: rotating fairness and the v4 baseline}
\code{scan} starts at a rotating equation pointer and takes the first eligible
equation and its cursor-selected solution. It performs no ranking or geometric
calculation. Individual cursor checks are constant work, and covered entries are
amortized over cursor advancement. Choose \code{scan} for the closest direct v4
queue comparison, debugging, or minimum dispatch overhead.

\subsection{\code{bandit}: learned productivity (default)}
Every equation carries a discounted discovery gain
\[
g_e\leftarrow0.7g_e+0.3d_e,\qquad g_e^{(0)}=5,
\]
where $d_e$ is the number of new solutions returned by its latest completed
trace. The optimistic initial value promotes exploration. Each dispatch sorts
equations by decreasing gain and takes the first eligible one, costing
$O(n_{\rm eq}\log n_{\rm eq})$ with $O(n_{\rm eq})$ gain storage.
After $2n_{\rm eq}$ consecutive no-discovery traces, the run switches one-way
to \code{scan} for the mechanical coverage tail. A discovery before that
threshold resets the counter.

\subsection{\code{novelty}: gain ranking plus geometric separation}
\code{novelty} uses the bandit equation rank, then chooses a separated start on
that equation. It samples at most 96 eligible candidates and 96 covered
references, maps them through a fixed eight-dimensional random projection,
includes both reference signs, and chooses the candidate with greatest
nearest-reference squared distance. Sampling is capped at 400 draws per
dispatch. With caps $k_c$, $k_r$, and projection dimension $k_p$, distance work
is bounded by $O(k_c k_r k_p)$. Dedicated seeded streams isolate this calculation
from unrelated RNG calls, but asynchronous completion remains nondeterministic.

\subsection{Selecting and choosing a policy}
\begin{lstlisting}[style=mat]
global HEBCPOLICY
HEBCPOLICY = 'scan';       % rotating v4-oriented baseline
HEBCPOLICY = 'bandit';     % default: discounted equation gain
HEBCPOLICY = 'novelty';    % bandit plus projected separation
[result,solutions] = run_merged_case('case57mod','parfeval',false);
result.scheduler
result.trace_metrics.trace_to_90pct
\end{lstlisting}
Unset or empty selects \code{bandit}. The legacy internal name \code{diverse} is
a deprecated alias for \code{bandit}; unknown names raise an error.
\code{trace_policy_config.m} normalizes the value through every entry point, so
\code{main.m} cannot silently overwrite an explicit choice.

\begin{center}\small
\begin{tabular}{@{}p{0.16\textwidth}p{0.25\textwidth}p{0.27\textwidth}p{0.22\textwidth}@{}}
\toprule
Policy & Objective & Cost and behavior & Best use\\
\midrule
\code{scan} & Fair coverage & Lowest selection cost & v4 comparison/debugging\\
\code{bandit} & Productive equations first & Gain sort; cheap tail hand-off &
general use; default\\
\code{novelty} & Separated starts first & Bounded projection-distance work &
early/anytime discovery\\
\bottomrule
\end{tabular}
\end{center}

Scheduler choice changes traversal order, not the power-flow model. Correctness
is checked by final count, residual, and order-independent set distance. Wall
time and trace-to-target are measured quantities, not guarantees. A few extra
traces can complete when already-running futures finish after another worker
establishes coverage; this benign tail redundancy explains occasional
one-trace differences between policies.

\subsection{Solution-connectivity diagrams}
Within each case all panels use identical graph geometry. Nodes are solutions,
color is discovery order, and edges join solutions covered by a shared trace.
Panel differences therefore show traversal history, not different accepted
final sets; the operating-point seed is black.

\IfFileExists{../connectivity_case14mod_scan.png}{
\begin{figure}[htbp]\centering
\begin{tabular}{ccc}
\includegraphics[width=0.30\textwidth]{../connectivity_case14mod_scan.png}&
\includegraphics[width=0.30\textwidth]{../connectivity_case14mod_bandit.png}&
\includegraphics[width=0.30\textwidth]{../connectivity_case14mod_novelty.png}\\
\code{scan}&\code{bandit}&\code{novelty}
\end{tabular}
\caption{case14mod connectivity for the common 30-solution set.}
\end{figure}}{}

\IfFileExists{../connectivity_case39_scan.png}{
\begin{figure}[htbp]\centering
\begin{tabular}{ccc}
\includegraphics[width=0.30\textwidth]{../connectivity_case39_scan.png}&
\includegraphics[width=0.30\textwidth]{../connectivity_case39_bandit.png}&
\includegraphics[width=0.30\textwidth]{../connectivity_case39_novelty.png}\\
\code{scan}&\code{bandit}&\code{novelty}
\end{tabular}
\caption{case39 connectivity for the validated common 176-solution set.}
\end{figure}}{}

A complete single-session \code{solution_connectivity.m} graph should contain
one connected component. Multiple components can indicate merged or inconsistent
bookkeeping rather than superior scheduler exploration.
